\subsection{The \texttt{rates} package}\label{subsec:appC-rates}

The multi-curve bootstrap and the time-changed SABR caplet pricer of Chapter~\ref{ch:prelim} are implemented as a small Python package, \texttt{thesis/code/rates/}, with a test file \texttt{thesis/code/tests/} \texttt{test\_rates.py}. The package depends on \texttt{numpy} only; the normal distribution function is taken from the standard library (\texttt{math.erf}), so no \texttt{scipy} installation is needed. The code is a reference implementation of the formulas of Chapter~\ref{ch:prelim}, not a production library: it has three modules and no calibration routine.

\paragraph{\texttt{curve.py}: discount curve and forward compounded rate.}
\texttt{bootstrap\_ois} returns discount factors $P(0,t_i)$ at increasing pillar maturities $t_1<\dots<t_n$ from the par rates $R_i$ of overnight index swaps maturing at $t_i$. The floating leg of an OIS is worth $1-P(0,t_i)$ at inception; the fixed leg is taken to pay $R_i(t_j-t_{j-1})$ at every pillar $t_j\le t_i$, so that the bootstrap is the closed recursion
\[
P(0,t_i)=\frac{1-R_i\sum_{j<i}(t_j-t_{j-1})P(0,t_j)}{1+R_i\,(t_i-t_{i-1})},
\]
with no root finding. The convention that fixed coupons fall on the pillar grid is exact when the pillars beyond one year are annual, which is the case for the curves used in Chapter~\ref{ch:usd}; a general coupon schedule would replace the annuity sum and nothing else. \texttt{discount} interpolates $\log P(0,t)$ linearly in $t$ between pillars, that is, with a constant continuously compounded forward rate on each interval, with $P(0,0)=1$ at the short end and a flat forward beyond the last pillar. \texttt{forward\_compounded} returns $F_0=(P(0,T)/P(0,T+\Delta)-1)/\Delta$, the identity $1+\Delta F_0=P(0,T)/P(0,T+\Delta)$ of Section~\ref{sec:prelim-compounded} under the continuous approximation.

\paragraph{\texttt{clock.py}: decay profile and clock.}
\texttt{g\_linear} is the profile of Assumption~\ref{ass:prelim-linear}, $g(t)=1$ for $t\le T$ and $(T+\Delta-t)/\Delta$ on $(T,T+\Delta]$. \texttt{tau} computes the clock $\tau(t)=\int_0^tg(u)^2\,du$ by composite Simpson quadrature on $[T,t]$, adding $T$ for the part where $g=1$. The value $\tau(T+\Delta)=T+\Delta/3$ is deliberately not hard-coded: it is a consequence of the linear-decay assumption, and the test verifies it numerically to $10^{-12}$ (Simpson's rule is exact for the quadratic $g^2$, so the agreement is to rounding). Any other profile can be passed as \texttt{g}; the pricer then uses its clock, and the factor $\Delta/3$ is replaced by whatever $c_g$ the profile implies, as in Example~\ref{ex:prelim-cg}.

\paragraph{\texttt{sabr.py}: Hagan volatility, Bachelier price, two caplet variants.}
\texttt{hagan\_normal\_vol} is formula \eqref{eq:prelim-hagan-N}, reducing to \eqref{eq:prelim-hagan-N0} at $\beta=0$; at $\zeta=0$ it uses $\zeta/x(\zeta)=1$. \texttt{bachelier\_call} is the undiscounted Bachelier call of Definition~\ref{def:prelim-implied}. \texttt{caplet\_time\_changed} is Lemma~\ref{lem:prelim-tc} in code: the Hagan volatility at the clock $\tau(T+\Delta)$, the Bachelier price at that clock, times $P(0,T+\Delta)$, per unit accrual. Both the compounded and the JIBAR caplet are priced under $\Q^{T+\Delta}$, and the code contains no convexity adjustment. \texttt{caplet\_nondecaying\_volvol} is the variant of Proposition~\ref{prop:prelim-novol} and is labelled in the source as an approximation, because no closed form exists for time-dependent vol-of-vol. Its two accountings are the two of the gap box after Proposition~\ref{prop:prelim-nueff}: \texttt{accounting="total"} uses the total-variance identification $\nu_{\mathrm{eff}}^2=\nu^2(T+\Delta)/\tau^*$ with $\rho$ unchanged; \texttt{accounting="hlw"} uses the effective parameters \eqref{eq:prelim-eff} computed by \texttt{hlw\_effective}. The latter evaluates the two integrals of \eqref{eq:prelim-eff} in calendar time by the substitution used in the derivation of Proposition~\ref{prop:prelim-novol}, which removes the $w^{-2/3}$ singularity of \eqref{eq:prelim-nutilde-linear} so that Simpson quadrature suffices.

\paragraph{Tests and what they reproduce.}
The test file runs under \texttt{pytest} or as a plain script. It checks: that the bootstrap returns $P(0,n)=(1+r)^{-n}$ on a flat annual curve and reprices its inputs to $10^{-14}$; that log-linear interpolation keeps a flat curve flat between and beyond the pillars; that $\tau(T+\Delta)=T+\Delta/3$ to $10^{-12}$ for two $(T,\Delta)$ pairs and that $\tau(t)=T+\tfrac\Delta3(1-g(t)^3)$ inside the window; that the at-the-money Bachelier price equals $\sigma\sqrt t/\sqrt{2\pi}$ to $10^{-16}$ and the deep-in and deep-out-of-the-money limits are the intrinsic value and zero; and that every figure of the worked example of Section~\ref{subsec:prelim-worked} is reproduced to the five significant figures printed there, from the clock $\tau^*$ through the at-the-money and $\mp50$ basis-point volatilities at the clock and quoted at $T+\Delta$ to the caplet price and the $7.1\%$ reduction against constant-parameter SABR without decay. The two vol-of-vol accountings are asserted against the same section: the differences from Lemma~\ref{lem:prelim-tc} at $\mp50$ and at the money are tenths of a basis point on the total-variance accounting and below $0.003$ basis points on the effective-parameter accounting, the values the double-precision run gives being $[+0.418,\,+0.125,\,+0.147]$ and $[+0.0023,\,-0.0000,\,-0.0023]$ basis points. The test also asserts the closed forms $\nu_{\mathrm{eff}}^2/\nu^2=1+\tfrac{2}{567}\Delta^3/(\tau^{*3}/3)$ and $\rho_{\mathrm{eff}}\nu_{\mathrm{eff}}/(\rho\nu)=(T\tau^*-T^2/2+\Delta^2/15)/(\tau^{*2}/2)$ to $10^{-9}$, which fixes the quadrature as well as the formulas.

\datanote{The test asserts the closed forms quoted in the gap box after Proposition~\ref{prop:prelim-nueff}, and the effective-parameter figures above are the values the code returns for them. Run: \texttt{cd thesis/code \&\& python3 tests/test\_rates.py} (or \texttt{python3 -m pytest tests/}); six tests, all passing on 9~September 2026 with Python~3.14 and numpy~2.4 and with Python~3.12 under pytest.}

\subsection{The \texttt{certify} package}\label{subsec:appC-certify}

The certificate of Chapter~\ref{ch:certification} is implemented as a second package, \texttt{thesis/code/} \texttt{certify/}: one module, a test file \texttt{certify/tests/test\_certify.py} and a figure script \texttt{certify/} \texttt{fig\_example.py}. It depends on \texttt{numpy} and, for the Beta and normal quantile functions, on \texttt{scipy}; because no system Python used in this work carries \texttt{scipy}, the run command is \texttt{cd thesis/code \&\& uv run --with numpy --with scipy python certify/tests/test\_certify.py} (seven tests, about fifteen seconds). The package regenerates Table~\ref{tab:example} and the series of the block-length sweep figure of Chapter~\ref{ch:sharpness} of Chapter~\ref{ch:certification}.

\paragraph{Blocking.}
\texttt{blocks} lays out the record of Section~\ref{sec:cert-scores} in the pattern of Algorithm~\ref{alg:prelim-blocking}: $n=\lfloor(N-s_0)/(2b)\rfloor$ blocks of length $b$ separated by gaps of length $b$, with the trailing gap held inside the record, so that with no fitting prefix $n=\lfloor N/(2b)\rfloor$. \texttt{deployment\_start} returns the first index of the gap-separated deployment block, and \texttt{block\_scores} sums the signed per-period hedging error over each block, which is the score of Definition~\ref{def:score}, for a single path or an array of paths. The block count is the one Chapter~\ref{ch:certification} uses ($250+240+10=500$ at $N=500$, $b=10$); a count that allows a block to open inside the trailing gap gives thirteen blocks at $b=20$ against the twelve of Table~\ref{tab:example}, and is wrong for that reason.

\paragraph{Threshold and certificate.}
\texttt{conformal\_k} and \texttt{conformal\_threshold} are $k=\lceil(1-\alpha)(n+1)\rceil$ and $\hat q=S_{(k)}$ of Definition~\ref{def:prelim-split}, returning $+\infty$ when $k>n$. \texttt{certificate} returns both forms of Theorem~\ref{thm:cov}: the marginal level $1-\alpha-\dK-(n+1)\beta(b)$ and the calibration-conditional level $b_{n,k}(\delta)-\dK-\beta(b)/\delta'$ with its confidence $1-\delta-n\beta(b)$. Its parts are exposed separately: \texttt{beta\_quantile} is the exact $\Beta(k,n+1-k)$ quantile, \texttt{dkw\_level} the weaker form of Corollary~\ref{cor:prelim-beta-bound}, \texttt{coupling\_cost} the $(n+1)\beta(b)$ of Lemma~\ref{lem:firstblock}. Two invariants of the method are enforced by the interface rather than left to documentation. The mixing coefficient enters as a function argument, \texttt{beta\_fn}, and is never estimated from the sample, because Proposition~\ref{prop:noest} shows that it cannot be. The drift term $\dK$ likewise enters as an argument, either a stress value or the plug-in bound below, because the certificate never sees a deployment score. Read as a sequence of calls, the function is Algorithm~\ref{alg:cert}.

\paragraph{The drift term.}
\texttt{ecdf\_gap} and \texttt{dK\_hat} give the two-sample Kolmogorov distance on the pooled grid; \texttt{W1\_hat} is the empirical $\Wone$ computed as the area between the two distribution functions, which is Lemma~\ref{thm:prelim-w1-identity} in code; \texttt{sqrt\_drift\_bound} is the bound $\sqrt{2L\Wone}$ of Proposition~\ref{prop:w1}. The linear form $L\Wone$ is not implemented, and the test exhibits the pair $\delta_0$ against $U[0,1/L]$ on which it fails.

\paragraph{The blocking curve.}
\texttt{deficit\_curve} evaluates $f(b)=\sqrt{b/N}+Nb^{-(r+1)}$, the expected calibration-conditional shortfall bound of Theorem~\ref{thm:cov}(ii) over the mixing class $\Bclass{r}$, which Chapter~\ref{ch:sharpness} minimises. \texttt{optimal\_block\_length} returns the exact minimiser $b^*=(2(r+1))^{2/(2r+3)}N^{3/(2r+3)}$ and \texttt{optimal\_deficit} its value $f(b^*)=\frac{2r+3}{2(r+1)}(2(r+1))^{1/(2r+3)}N^{-r/(2r+3)}$. The bare balance $N^{3/(2r+3)}$, at which the two terms are equal, is kept under the separate name \texttt{balance\_block\_length} and is not the minimiser: at $N=500$, $r=2$ the minimiser is $23.93$ and the balance $14.35$, a factor $1.67$. No function in the package returns the balance under a name that could be read as optimal. \texttt{expected\_shortfall\_bound} evaluates the same object with the actual constants, $1/(2\sqrt{n+2})+(n+1)\beta(b)$, for a supplied $\beta$; for the AR(1) of Chapter~\ref{ch:certification} with $\varphi=\tfrac12$ it is smallest at $b=11$ ($0.1077$, against $0.1089$ at $b=10$), which is a property of a geometrically mixing process rather than an instance of $b^*$, since geometric mixing lies in every $\Bclass{r}$ and $b^*$ depends on $r$.

\paragraph{Comparison, not theorem.}
\texttt{stationary\_bootstrap\_quantiles} is the stationary bootstrap of \citet{politis1994} applied to the $(1-\alpha)$ quantile of the overlapping block sums. It is an empirical comparator and carries no coverage statement. On the AR(1) example without drift ($200$ replications, $200$ bootstrap resamples) the conformal threshold covers the deployment score with frequency $0.935$, the median bootstrap quantile with $0.860$ and its $95\%$ upper envelope with $0.910$; the mean thresholds are $7.70$, $6.00$ and $7.61$ against a true $P$-quantile of $6.54$. The bootstrap estimates a quantile of the calibration law and under-covers a future block, which is why the thesis states no theorem about it.

\paragraph{The AR(1) example.}
\texttt{ar1\_beta\_bound} takes $\varphi$ and returns the map $k\mapsto\tfrac12\sqrt{-\log(1-\varphi^{2k})}$, the first inequality of Lemma~\ref{prop:prelim-ar1} and the bound Table~\ref{tab:example} uses, in the form \texttt{certificate} expects for \texttt{beta\_fn}; \texttt{ar1\_beta\_exact} evaluates $\beta(k)$ by quadrature, giving $3.108\times10^{-4}$ at $\varphi=\tfrac12$, $k=10$, so the bound is a factor $1.57$ loose. \texttt{ar1\_block\_var} and \texttt{ar1\_block\_sd} return the block variance $v_b$, \texttt{ar1\_drift\_dK} the drift distance $\dK=2\Phi\big(b\mu/(2\sqrt{v_b})\big)-1$, and \texttt{simulate\_ar1} draws stationary paths from a supplied innovation standard deviation, to which the caller passes $\sqrt{1-\varphi^2}$ so that the stationary variance is one; the normalisation is stated with every figure, because it changes every threshold while leaving every coverage number unchanged.

\paragraph{Tests and what they reproduce.}
The seven tests are: the blocking layout and the block count above; that \texttt{ar1\_beta\_bound} dominates \texttt{ar1\_beta\_exact} on a grid of $\varphi$ and $k$; \texttt{test\_table\_3\_1}, which reproduces every cell of Table~\ref{tab:example} at $b\in\{5,10,20\}$ with $\mu=0.05$; a Monte Carlo of the whole procedure at $b=10$ ($10^4$ replications, seed $20260909$) giving realised coverage $0.9113$ with standard error $0.0028$, mean $F_Q(\hat q)=0.9096$ against the coupled value $0.9089$, and $\Prob(F_Q(\hat q)\ge0.7750)=0.9710$ against the certified confidence $0.9378$; the drift bound of Proposition~\ref{prop:w1} on a Gaussian pair, where $\dK=0.1974$, $\hat\dK=0.2005$, $\hat\Wone=2.952$ and $\sqrt{2L\hat\Wone}=0.632$; the minimiser and the balance above; and the bootstrap comparison. \texttt{fig\_example.py} sweeps $b=2,\dots,50$ at $\varphi\in\{0.5,0.8\}$ with $10^4$ replications per point and the same seed and writes \texttt{certify/fig\_example.csv}, the series of the block-length sweep figure of Chapter~\ref{ch:sharpness}: breach frequencies $0.109$, $0.093$, $0.098$ at $b=5,10,20$ for $\varphi=\tfrac12$ and $0.112$, $0.089$, $0.091$ for $\varphi=0.8$. Each lies below the certified breach bound $\alpha+(n+1)\beta(b)+(1-k/(n+1))+\dK$, whose excess over $\alpha$ is $0.925$, $0.129$ and $0.130$ at $\varphi=\tfrac12$ and $8.72$, $1.50$ and $0.186$ at $\varphi=0.8$; almost all the slack is in the coupling term $(n+1)\beta(b)$, which is why Chapter~\ref{ch:sharpness} minimises over $b$ rather than reporting a single block length.

\datanote{Table~\ref{tab:example} and the block-length sweep figure of Chapter~\ref{ch:sharpness} of Chapter~\ref{ch:certification} are regenerated by this package, not computed by hand: \texttt{test\_table\_3\_1} and \texttt{fig\_example.py} at seed $20260909$. The regenerated cells are $\beta(b)=1.563\times10^{-2}$, $4.883\times10^{-4}$, $4.768\times10^{-7}$; $(n+1)\beta(b)=0.7971$, $0.0127$, $6.2\times10^{-6}$; $\dK=0.0299$, $0.0391$, $0.0533$; marginal level $0.0730$, $0.8482$, $0.8467$; $b_{n,k}(0.05)=0.8262$, $0.8239$, $0.7791$; $\beta(b)/\delta'=0.3126$, $0.0098$, $9.5\times10^{-6}$; conditional level $0.4837$, $0.7750$, $0.7258$; confidence $0.1686$, $0.9378$, $0.9500$; and $\E F_Q(\hat q^*)=0.90889$ by quadrature against the $\Beta(24,2)$ density. Every one agrees with the digits printed in Chapter~\ref{ch:certification}, so no hand-arithmetic discrepancy remains to be recorded. The three valid bounds on $\beta(10)$ at $\varphi=\tfrac12$ that appear in the thesis, $4.883\times10^{-4}$ from Lemma~\ref{prop:prelim-ar1}, $5.638\times10^{-4}$ from the $k$-uniform form $|\varphi|^k/(2\sqrt{1-\varphi^2})$ and $6.905\times10^{-4}$ from Proposition~\ref{prop:prelim-ar1}, are not corrected into one another; the code implements the first because that is the one the table uses.}

\gap{Three modules of the reproducibility package are not written. \texttt{strip} (caplet-volatility stripping from cap and floor quotes), \texttt{calibrate} (SABR calibration to a quoted smile with the parameter count of Chapter~\ref{ch:identifiability}) and \texttt{scores} (realised hedging-error scores of a fixed policy) are scheduled for programme months 5--9 and 12--17; nothing in this appendix or in Chapters~\ref{ch:usd} and~\ref{ch:zar} may be read as describing code that exists. Within \texttt{certify}, three limitations are recorded: $\dK$ is an input and is supplied by Chapter~\ref{ch:usd} from observed deployment scores with an assumed density bound $L$; the dependent-data rate for $\hat\Wone$ is not implemented, only the estimator; and \texttt{blocks} discards the tail of the record shorter than $2b$ rather than fitting one further block.}

\figplan{Module diagram, three columns. Left column, implemented: \texttt{rates} (\texttt{curve}, \texttt{clock}, \texttt{sabr}) and \texttt{certify} (blocking, threshold, certificate, drift, blocking curve, AR(1) helpers, bootstrap comparator). Middle column, the thesis objects each module computes, with arrows from the left: bootstrap and forward compounded rate, clock and time-changed caplet, Hagan volatility (Chapter~\ref{ch:prelim}); score, threshold, both forms of the certificate, drift bound, optimal block length (Chapters~\ref{ch:certification} and~\ref{ch:sharpness}). Right column, dashed boxes for the modules not yet written, \texttt{strip}, \texttt{calibrate} and \texttt{scores}, with dashed arrows to the caplet-volatility stripping, SABR calibration and hedging-error scores of Chapters~\ref{ch:usd} and~\ref{ch:zar}, annotated with the programme months in which they are scheduled.}{Code modules and the thesis objects they compute. Solid boxes are implemented and tested at the date of this appendix; dashed boxes are scheduled. The figure records which results of the thesis are reproducible from the package as it stands and which depend on code still to be written. Source: the package layout of Section~\ref{sec:appC-code} and the protocol of Section~\ref{sec:appC-protocol}.\label{fig:appC-modules}}
